\subsection{Semisimplicity}
\setcounter{subsection}{1}

\begin{exr}[title=Lemma 2.3] %NOTE
	Let $R$ be a ring and let $M$ be a nonzero $R$-module.

	Show that if $M$ is semisimple, then it must contain a simple submodule.
\end{exr} \

\begin{pf}
	asdf
\end{pf} \

\begin{exr}[title=Theorem 2.4] %NOTE
	Prove that the following are equivalent. \vspace{-0.2in}
	\begin{enumerate}[i)]
		\item $M$ is semisimple.
		\item $M$ can be expressed as a direct sum of simple modules,
			\textit{i.e.} $M=\bigoplus_{i\in I} M_i$ for some (possibly infinite)
			family of simple submodules $M_i$.
		\item $M$ can be expressed as a sum of simple modules,
			\textit{i.e.} $M=\sum_{i\in I} M_i$ for some (possibly infinite)
			family of simple submodules $M_i$.
	\end{enumerate}
\end{exr} \

\begin{pf}
	asdf
\end{pf} \

\begin{exr}
	Show that if $M$ is semisimple, then $M$ is finitely generated if and only if
	$M$ is a direct sum of \textit{finitely many} simple modules.
\end{exr} \

\begin{pf}
	Let $M=\bigoplus_{i}M_{i}$ be semisimple.

	If $M$ is finitely generated, say by $m_{1},\dots,m_{n}$,
	then $\pi_{i}(m_{k})=0$ for all but finitely many $i$'s; that is,
	$m_{k}$ only has finitely many non-zero coordinates.
	Since this holds for each $k$, and there are finitely many generators,
	then the direct sum must necessarily be finite.

	On the other hand, if $M=\bigoplus_{i=1}^{n}M_{i}$ is semisimple,
	fix some non-zero $m_{i}\in M_{i}$ for each $i$.
	Then, since $M_{i}$ is simple, $Mm_{i}=M_{i}m_{i}=M_{i}$, so
	the $m_{i}$ form a finite generating set for $M$.
\end{pf}


\nq
\begin{exr}[title=Theorem 2.5] %NOTE
	Let $R$ be a ring. Prove that the following are equivalent.
	(A \textbf{left semisimple ring} is one which satisfies one, hence all, of these conditions.) \vspace{-0.2in}
	\begin{enumerate}[i)]
		\item All short exact sequences of left $R$-modules must split.
			(Equivalently: if $M$ is a left $R$-module and
			$N$ is a submodule of $M$, then there exists a submodule $N'$ of $M$ such that
			$M=N\oplus N'$.)
		\item Every left $R$-module is semisimple.
		\item Every finitely generated left $R$-module is semisimple.
		\item All cyclic left modules are semisimple.
		\item $R$ is semisimple as a left module over itself.
	\end{enumerate}
\end{exr} \

\begin{pf}
	asdf
\end{pf}


\nq
\begin{exr}[title=Theorem 2.8] %NOTE
	Let $R$ be a ring. Prove that the following are equivalent.
	\begin{enumerate}[i)]
		\item $R$ is left semisimple.
		\item Every left $R$-module is projective.
		\item Every finitely generated left $R$-module is projective.
		\item All cyclic left modules are projective.
	\end{enumerate}
\end{exr} \

\begin{pf}
	asdf
\end{pf}


\nq
\begin{exr}[title=Theorem 2.9] %NOTE
	Let $R$ be a ring. Prove that the following are equivalent.
	\begin{enumerate}[i)]
		\item $R$ is left semisimple.
		\item Every left $R$-module is injective.
		\item Every finitely generated left $R$-module is injective.
		\item All cyclic left modules are injective.
	\end{enumerate}
\end{exr} \

\begin{pf}
	asdf
\end{pf}


\nq
\begin{exr}[title=Exercise 2.1 part a)] %TODO
	If $R$ is left semisimple, does it follow that
	every subring of $R$ is left semisimple?
\end{exr} \

\begin{pf}
	asdf
\end{pf} \

\begin{exr}[title=Exercise 2.1 part b)] %TODO
	Can every ring be embedded as a subring of a left semisimple ring?
\end{exr} \

\begin{pf}
	asdf
\end{pf}


\nq
\begin{exr}[title=Exercise 2.2]
	Let $R=\prod_{i\in I}F_i$ be a product of fields.
	Prove that $R$ is semisimple if and only if the product is finite.
\end{exr} \

\begin{pf}
	If the product is finite, $R$ is evidently semisimple by definition.
	On the other hand, if $R$ is semisimple, then $_{R}R$ is a semisimple
	module.
	Evidently, the simple submodules of $R$ are precisely the $F_{i}$'s,
	and since $_{R}R$ is finitely generated (by the unit):
	\begin{equation*}
		R \ = \ F_{i_{1}} \oplus \cdots \oplus F_{i_{n}}
		\ \simeq \ F_{i_{1}} \times \cdots \times F_{i_{n}}
	\end{equation*}
	by 1c).
\end{pf}


\nq
\begin{exr}[title=Exercise 2.3]
	What are the semisimple $\bZ$-modules?
\end{exr} \

\begin{pf}
	Let $M$ be a semisimple $\bZ$-module.
	Recall that $\bZ$-modules can be identified with abelian groups.
	Considering $M=\bigoplus_{i}M_{i}$, each $M_{i}$ is a simple $\bZ$-module,
	which we can think of as a subgroup of $M$ with no proper nontrivial subgroups.

	Each $M_{i}$ must be generated by a single element;
	if not, suppose $x,y\in M_{i}$ where $x$ does not generate $y$.
	Then we can consider $\ang{x}\leq M_{i}$;
	this would be a proper non-trivial subgroup of $M_{i}$.

	Writing $M_{i}=\ang{x_{i}}$, it is clear that $x_{i}$ must have prime order.
	Therefore, all the semisimple $\bZ$-modules are of the form:
	\begin{equation*}
		M \ = \ \bigoplus_{p\in P}\bZ_{p}
	\end{equation*}
	where $P$ is some possibly repeating (possibly infinitely many times)
	collection of prime numbers.
\end{pf}


\nq
\begin{exr}[title=Exercise 2.5] %TODO
	Let $R$ be a left semisimple ring.
	Show that if $I$ is a right ideal and $J$ is a left ideal, then
	$IJ=I\cap J$. Prove that:
	\begin{gather*}
		I\cap (J+K)= (I\cap J)+(I\cap K) \\
        I+(J\cap K) = (I+J)\cap (I+K)
	\end{gather*}
\end{exr}
